% ============================================================================
% AERIS Paper - Section 6: Results and Analysis
% Target Journal: Applied Intelligence (APIN) - Springer
% Version: 2026-01-18 (Data Authenticity Corrected)
% ============================================================================

\section{Results and Analysis}
\label{sec:results}

This section presents comprehensive experimental results comparing AERIS
against baseline protocols. All data presented are actual measurements from
200 independent simulation runs per configuration. We organize results around
five key dimensions: baseline comparison, scalability, dynamic topology
robustness, statistical validation, and ablation analysis.

\subsection{Baseline Comparison}

Table~\ref{tab:baseline_comparison} presents the fundamental performance
comparison at 100 nodes.

\begin{table}[htbp]
\centering
\caption{Baseline Protocol Comparison (100 Nodes, 200 Rounds, Measured Data)}
\label{tab:baseline_comparison}
\begin{tabular}{@{}lcccc@{}}
\toprule
Protocol & Energy (mJ) & PDR (\%) & Execution Time (s) \\
\midrule
LEACH & 100.7 & 100.0 & 0.10 \\
\textbf{PEGASIS} & \textbf{41.9} & 100.0 & \textbf{0.02} \\
HEED & 87.3 & 99.98 & 0.12 \\
\textbf{AERIS} & 82.1 & \textbf{100.0} & 1.14 \\
\bottomrule
\end{tabular}
\end{table}

\subsubsection{Energy Consumption Analysis}

\textbf{Honest Assessment}: PEGASIS achieves the lowest energy consumption
(41.9mJ), approximately \textbf{49\% lower} than AERIS (82.1mJ). This
substantial difference stems from PEGASIS's chain-based aggregation where
each node transmits only to its immediate neighbor, minimizing transmission
distances.

AERIS achieves 18.5\% lower energy consumption than LEACH (82.1mJ vs 100.7mJ),
positioning it between the two classical approaches.

\textbf{Key Finding}: For energy-critical applications, PEGASIS remains
the optimal choice. AERIS provides energy improvement only relative to LEACH.

\subsubsection{PDR Analysis}

At 100 nodes, both AERIS, LEACH, and PEGASIS achieve 100\% PDR. The differences
emerge at larger scales, as analyzed in Section~\ref{subsec:scalability}.

\subsection{Scalability Analysis}
\label{subsec:scalability}

This is AERIS's primary advantage: \textbf{PDR stability at scale}.

\subsubsection{PDR at Scale}

Table~\ref{tab:scalability_pdr} presents PDR measurements across network scales.

\begin{table}[htbp]
\centering
\caption{PDR Scalability Comparison (\%, Measured Data)}
\label{tab:scalability_pdr}
\begin{tabular}{@{}lccccc@{}}
\toprule
Protocol & 50 Nodes & 100 Nodes & 200 Nodes & 300 Nodes & 500 Nodes \\
\midrule
LEACH & 100.0 & 100.0 & 99.71 & 99.30 & \textbf{98.68} \\
PEGASIS & 100.0 & 100.0 & 100.0 & 100.0 & 100.0 \\
HEED & 100.0 & 99.98 & 99.95 & 99.88 & 99.72 \\
\textbf{AERIS} & \textbf{100.0} & \textbf{100.0} & \textbf{100.0} & \textbf{100.0} & \textbf{100.0} \\
\bottomrule
\end{tabular}
\end{table}

\textbf{Critical Finding}:
\begin{itemize}
    \item \textbf{LEACH PDR degrades 1.32\%} at 500 nodes (100\% $\rightarrow$ 98.68\%)
    \item AERIS and PEGASIS both maintain 100\% PDR
    \item This difference is statistically significant ($p < 0.001$)
\end{itemize}

The LEACH PDR degradation results from increased transmission distances causing
packet loss in the direct cluster-head-to-base-station communication at scale.

\subsubsection{Energy Scalability}

Table~\ref{tab:scalability_energy} presents energy consumption across scales.

\begin{table}[htbp]
\centering
\caption{Energy Consumption Scalability (mJ, Measured Data)}
\label{tab:scalability_energy}
\begin{tabular}{@{}lcccc@{}}
\toprule
Protocol & 50 Nodes & 100 Nodes & 300 Nodes & 500 Nodes \\
\midrule
LEACH & 48.2 & 100.7 & 328.4 & 562.1 \\
\textbf{PEGASIS} & \textbf{19.8} & \textbf{41.9} & \textbf{134.5} & \textbf{225.8} \\
HEED & 42.1 & 87.3 & 278.2 & 471.3 \\
AERIS & 39.5 & 82.1 & 262.8 & 445.6 \\
\bottomrule
\end{tabular}
\end{table}

PEGASIS maintains its energy efficiency advantage across all scales,
consuming approximately 50\% less energy than AERIS consistently.

\subsection{Dynamic Topology Robustness}

This section evaluates protocol robustness under node churn and regional failure.

\subsubsection{Node Churn Experiment}

Table~\ref{tab:churn_results} presents PDR under varying node failure rates.

\begin{table}[htbp]
\centering
\caption{PDR Under Node Churn (\%, Measured Data)}
\label{tab:churn_results}
\begin{tabular}{@{}lcccc@{}}
\toprule
Protocol & 0\% & 10\% & 20\% & 30\% \\
\midrule
LEACH & 100.0 & 100.0 & 100.0 & 100.0 \\
PEGASIS & 100.0 & 100.0 & 100.0 & 100.0 \\
HEED & 99.98 & 99.97 & 99.96 & 99.95 \\
\textbf{AERIS} & \textbf{100.0} & \textbf{100.0} & \textbf{100.0} & \textbf{100.0} \\
\bottomrule
\end{tabular}
\end{table}

All protocols demonstrate robust performance under node churn, with AERIS,
LEACH, and PEGASIS maintaining 100\% PDR even at 30\% node failure.

\subsubsection{Regional Failure Experiment}

Table~\ref{tab:regional_results} presents PDR under concentrated regional
failures (affecting up to 40\% of nodes in a localized area).

\begin{table}[htbp]
\centering
\caption{PDR Under Regional Failure (40\% Nodes Failed, Measured Data)}
\label{tab:regional_results}
\begin{tabular}{@{}lcc@{}}
\toprule
Protocol & Baseline & 40\% Regional Failure \\
\midrule
LEACH & 100.0 & $\sim$99.997 \\
PEGASIS & 100.0 & 100.0 \\
HEED & 99.98 & 99.94 \\
\textbf{AERIS} & \textbf{100.0} & \textbf{100.0} \\
\bottomrule
\end{tabular}
\end{table}

AERIS maintains 100\% PDR even with 40\% regional node failure, demonstrating
strong spatial robustness through gateway redundancy mechanisms.

\subsection{Statistical Validation}

All reported results are validated using rigorous statistical methods.

\subsubsection{Significance Testing}

Table~\ref{tab:significance} reports Welch's t-test results for
AERIS vs. baselines at 500 nodes.

\begin{table}[htbp]
\centering
\caption{Statistical Significance (AERIS vs. Baselines, 500 Nodes)}
\label{tab:significance}
\begin{tabular}{@{}lcccc@{}}
\toprule
Comparison & Metric & Difference & $p$-value & Cohen's $d$ \\
\midrule
AERIS vs. LEACH & PDR & +1.32\% & $<$0.001*** & 1.89 (large) \\
AERIS vs. LEACH & Energy & -116.5mJ & $<$0.001*** & 2.76 (large) \\
AERIS vs. PEGASIS & Energy & +219.8mJ & $<$0.001*** & 4.15 (large) \\
AERIS vs. HEED & PDR & +0.28\% & 0.002** & 0.70 (medium) \\
\bottomrule
\end{tabular}
\begin{tablenotes}
\small
\item *** $p < 0.001$; ** $p < 0.01$; * $p < 0.05$ (Holm-Bonferroni corrected)
\end{tablenotes}
\end{table}

\subsubsection{Effect Size Interpretation}

\begin{itemize}
    \item \textbf{AERIS vs. LEACH (PDR)}: Cohen's $d = 1.89$ indicates
    AERIS's reliability advantage at scale is practically significant.

    \item \textbf{AERIS vs. PEGASIS (energy)}: Cohen's $d = 4.15$ confirms
    PEGASIS's substantial energy advantage is statistically and practically
    significant.
\end{itemize}

\subsection{Honest Trade-off Summary}

Table~\ref{tab:honest_summary} presents the honest comparison at 500 nodes,
highlighting each protocol's optimal domain.

\begin{table}[htbp]
\centering
\caption{Honest Trade-off Summary (500 Nodes, Measured Data)}
\label{tab:honest_summary}
\begin{tabular}{@{}lccccc@{}}
\toprule
Metric & AERIS & PEGASIS & LEACH & HEED & Winner \\
\midrule
\textbf{PDR} & \textbf{100\%} & \textbf{100\%} & 98.68\% & 99.72\% & AERIS/PEGASIS \\
\textbf{Energy} & 445.6mJ & \textbf{225.8mJ} & 562.1mJ & 471.3mJ & \textbf{PEGASIS} \\
\textbf{Exec. Time} & 15.17s & \textbf{0.11s} & 1.03s & 0.64s & \textbf{PEGASIS} \\
Scale Reliability & \checkmark & \checkmark & $\times$ & $\times$ & AERIS/PEGASIS \\
\bottomrule
\end{tabular}
\end{table}

\textbf{Conclusions}:
\begin{itemize}
    \item \textbf{AERIS vs. LEACH}: AERIS provides better PDR stability and
    lower energy consumption
    \item \textbf{PEGASIS vs. AERIS}: PEGASIS achieves lower energy and faster
    execution
    \item \textbf{AERIS value proposition}: Better reliability than LEACH with
    acceptable energy overhead
\end{itemize}

\subsection{Ablation Study}

Table~\ref{tab:ablation} presents the contribution of each AERIS component.

\begin{table}[htbp]
\centering
\caption{AERIS Ablation Study (100 Nodes, Measured Data)}
\label{tab:ablation}
\begin{tabular}{@{}lccc@{}}
\toprule
Configuration & PDR (\%) & Energy (mJ) & Hedges' $g$ \\
\midrule
AERIS Full & 100.0 & 82.1 & --- \\
$-$Gateway & 89.2 & 75.4 & \textbf{10.09} (large) \\
$-$Fairness & 99.4 & 84.3 & 0.59 (medium) \\
$-$Safety & 99.7 & 81.8 & 0.42 (small) \\
$-$CAS & 100.0 & 82.0 & \textbf{0.00} (negligible) \\
\bottomrule
\end{tabular}
\end{table}

\textbf{Key Finding}: The Gateway component is the primary contributor to
AERIS's reliability (Hedges' $g = 10.09$). The CAS module shows negligible
independent effect ($g = 0.00$), functioning as a supporting mechanism
rather than the core innovation.

\subsection{Data Authenticity Statement}

All results in this section are actual measurements from simulation experiments.
The following data sources were used:
\begin{itemize}
    \item \texttt{results/comprehensive\_dynamic\_experiments.json}
    \item 200 independent runs per configuration
    \item Same random seeds ensuring reproducibility
    \item Welch's t-test with Holm-Bonferroni correction
\end{itemize}
